\chapter{Mathematical Background}
\label{app:math}

\section{Bianchi's IEEE~802.11 Saturation Throughput Model}

Bianchi~\cite{bianchi2000} derived a closed-form expression for the saturation
throughput of IEEE~802.11 DCF. The model assumes:

\begin{itemize}
  \item $n$ nodes, each always having a packet to transmit (saturated conditions).
  \item Transmission probability $\tau$ per slot, derived from the steady-state analysis
    of the back-off Markov chain.
  \item No hidden terminals; all nodes within mutual range.
\end{itemize}

The key equations are:

\begin{equation}
  \tau = \frac{2(1 - 2p)}{(1 - 2p)(W + 1) + pW\!\left(1 - (2p)^m\right)}
  \label{eq:tau}
\end{equation}

where $p$ is the conditional collision probability, $W = \text{CW}_{\min}$, and $m$ is
the maximum back-off stage such that $2^m W = \text{CW}_{\max}$.

The collision probability is:

\begin{equation}
  p = 1 - (1 - \tau)^{n-1}
  \label{eq:collision_prob}
\end{equation}

These equations form a system that can be solved iteratively for $\tau$ and $p$. The
saturation throughput $S$ is then:

\begin{equation}
  S = \frac{P_s \cdot P_{tr} \cdot E[L]}
         {(1 - P_{tr})\sigma + P_{tr}P_s T_s + P_{tr}(1-P_s)T_c}
  \label{eq:saturation}
\end{equation}

where:
\begin{align}
  P_{tr} &= 1 - (1-\tau)^n & &\text{(probability at least one node transmits)} \\
  P_s    &= \frac{n\tau(1-\tau)^{n-1}}{P_{tr}} & &\text{(prob.\ of successful transmission)} \\
  T_s    &= \text{DIFS} + T_{\text{frame}} + \text{SIFS} + T_{\text{ACK}} & &\text{(successful slot time)} \\
  T_c    &= \text{DIFS} + T_{\text{frame}} & &\text{(collision slot time)}
\end{align}

For 802.11g at 54~Mbps with 1000~B payload, $T_{\text{frame}} \approx 200\,\mu\text{s}$.
As $n$ increases from 10 to 50, $P_s$ decreases and $T_c$ events become more frequent,
causing $S$ to decrease from near-maximum to well below 50\% of the theoretical channel
capacity. This is the mechanism that explains why DSDV's additional control-frame
traffic causes disproportionate throughput degradation at $N = 50$.

\section{Retry-Limit Frame Drop Probability}

When the channel is saturated and the collision probability $p$ is high, the probability
that a frame survives all $m_r$ retry attempts without an abort is:

\begin{equation}
  P_{\text{delivery}} = (1 - p)^{m_r + 1} = (1-p)^8 \qquad (m_r = 7)
  \label{eq:drop_prob}
\end{equation}

Table~\ref{tab:drop_prob} shows how frame-drop probability varies with collision
probability:

\begin{table}[ht]
\centering
\caption{Frame-drop probability vs.\ collision probability ($m_r = 7$ retries).}
\label{tab:drop_prob}
\begin{tabular}{@{}ccc@{}}
\toprule
\textbf{Collision probability $p$} & \textbf{Delivery probability} & \textbf{Drop rate} \\
\midrule
0.01 & 0.923 & 7.7\% \\
0.05 & 0.663 & 33.7\% \\
0.10 & 0.430 & 57.0\% \\
0.20 & 0.168 & 83.2\% \\
0.30 & 0.058 & 94.2\% \\
0.50 & 0.004 & 99.6\% \\
\bottomrule
\end{tabular}
\end{table}

This table demonstrates that even a modest increase in collision probability --- from 5\%
to 10\% --- nearly doubles the drop rate. DSDV's routing broadcasts, by increasing the
effective number of competing nodes, can push the network from the 5\% to the 10--20\%
collision-probability regime, causing the catastrophic drop rates observed at $N = 50$.
